\hypertarget{appendix-40df0855b4a7}{}
\section*{Automatic differentiation and backpropagation}
\chapter{Automatic differentiation and backpropagation}\label{appendix-40df0855b4a7-automatic-differentiation-and-backpropagation}

\textbf{Curriculum depth:} Developed · \textbf{First seen:} University courses · \textbf{Latest development:} Phase 9

\section{Definition and mechanism}\label{appendix-40df0855b4a7-definition-and-mechanism}

Automatic differentiation evaluates derivatives by applying the chain rule to the operations in a program. Each operation has a local derivative rule; composing the rules gives the derivative of the complete computation without symbolically expanding it or approximating it by finite differences. Forward mode propagates a tangent perturbation and computes a Jacobian-vector product (JVP), useful when there are few inputs. Reverse mode propagates sensitivities from outputs back to inputs and computes a vector-Jacobian product (VJP), making gradients of one scalar loss with respect to millions of parameters practical. Backpropagation is reverse-mode differentiation through a neural network's computation graph. It does not remove numerical problems: nondifferentiable operations require conventions, saturated activations can yield tiny gradients, and an incorrect graph or reduction yields an incorrect gradient. Gradients are local sensitivity information, not proof of a global optimum. Finite differences remain a useful small-scale check, but cancellation and step-size tradeoffs make them a poor training method.

\subsection{Phase 9--10 development}\label{appendix-40df0855b4a7-phase-910-development}

Forward mode propagates an input tangent through local derivatives and computes a Jacobian-vector product. Reverse mode propagates an output sensitivity backward and computes a vector-Jacobian product. These are chain-rule evaluations on the operation graph, distinct from finite differences, which perturb inputs and can suffer cancellation. The reverse pass requires a valid reverse-topological order and storage that remains alive until propagation finishes. Differentiating a converged solution uses an implicit equation and its local Jacobian; it is not automatically the same as differentiating every solver iteration.

\section{Prerequisites}\label{appendix-40df0855b4a7-prerequisites}

\begin{itemize}
\tightlist
\item
  \href{../04-Visual-Computing-II/main.pdf\#nameddest=appendix-29f4bea4c2bb}{Constrained optimization and inverse kinematics (04-Visual-Computing-II, appendix-29f4bea4c2bb)}
\item
  \href{../02-Foundations-II/main.pdf\#nameddest=appendix-ab31ef447dbf}{Matrix multiplication and composition order (02-Foundations-II, appendix-ab31ef447dbf)}
\item
  \href{../01-Foundations-I/main.pdf\#nameddest=appendix-c2ac4a84c283}{Floating-point reassociation (01-Foundations-I, appendix-c2ac4a84c283)}
\item
  \href{../02-Foundations-II/main.pdf\#nameddest=appendix-7bc181935c42}{Derivative as a local linear map (02-Foundations-II, appendix-7bc181935c42)}
\item
  \href{../01-Foundations-I/main.pdf\#nameddest=appendix-d8a7ddb675eb}{Jacobians, pushforwards, and pullbacks (01-Foundations-I, appendix-d8a7ddb675eb)}
\end{itemize}

\section{Used by}\label{appendix-40df0855b4a7-used-by}

\begin{itemize}
\tightlist
\item
  \href{Continuous\%20Notes/University\%20-\%20Worlds,\%20Images\%20and\%20Inference\%20-\%20Graphics,\%20Vision\%20and\%20Photography\#Learning\%20a\%20visual\%20classifier\%20from\%20examples}{Learning a visual classifier from examples}
\item
  \href{Continuous\%20Notes/University\%20-\%20Worlds,\%20Images\%20and\%20Inference\%20-\%20Graphics,\%20Vision\%20and\%20Photography\#Optimization,\%20constraints,\%20and\%20inverse\%20kinematics}{Optimization, constraints, and inverse kinematics}
\item
  \href{../04-Visual-Computing-II/main.pdf\#nameddest=appendix-47c51ea026fb}{Learned image features and CNNs (04-Visual-Computing-II, appendix-47c51ea026fb)}
\item
  \href{../04-Visual-Computing-II/main.pdf\#nameddest=appendix-a9d5dfdd164a}{Attention and visual transformers (04-Visual-Computing-II, appendix-a9d5dfdd164a)}
\item
  \href{../04-Visual-Computing-II/main.pdf\#nameddest=appendix-2c568f05ee34}{Latent-variable models and generative flows (04-Visual-Computing-II, appendix-2c568f05ee34)}
\item
  \href{../04-Visual-Computing-II/main.pdf\#nameddest=appendix-c5b332e400e7}{Radiance fields and Gaussian splatting (04-Visual-Computing-II, appendix-c5b332e400e7)}
\item
  \href{../04-Visual-Computing-II/main.pdf\#nameddest=appendix-b0c42da325f1}{Perception feedback and robot learning (04-Visual-Computing-II, appendix-b0c42da325f1)}
\item
  \href{../02-Foundations-II/main.pdf\#nameddest=appendix-8836553d56fc}{Implicit differentiation and parameter sensitivities (02-Foundations-II, appendix-8836553d56fc)}
\end{itemize}

\section{Full teaching explanation}\label{appendix-40df0855b4a7-full-teaching-explanation}

\begin{itemize}
\tightlist
\item
  \href{Continuous\%20Notes/University\%20-\%20Worlds,\%20Images\%20and\%20Inference\%20-\%20Graphics,\%20Vision\%20and\%20Photography\#Learning\%20a\%20visual\%20classifier\%20from\%20examples}{Learning a visual classifier from examples}
\item
  \href{Continuous\%20Notes/University\%20-\%20Worlds,\%20Images\%20and\%20Inference\%20-\%20Graphics,\%20Vision\%20and\%20Photography\#Optimization,\%20constraints,\%20and\%20inverse\%20kinematics}{Optimization, constraints, and inverse kinematics}
\item
  \href{Continuous\%20Notes/Computing\%20Foundations\%20-\%20From\%20Bits\%20to\%20C++\%20and\%20Systems\%20APIs\#Numerical\%20differentiation\%20and\%20automatic\%20differentiation}{Phase 9 teaching section}
\item
  \href{../01-Foundations-I/main.pdf\#nameddest=appendix-adf48f5dc5c1}{Visual Recall/09 - Calculus and Optimization - Visual Recall.md (01-Foundations-I, appendix-adf48f5dc5c1)} · Phase 9 visual recall
\end{itemize}

\section{Recall and next depth}\label{appendix-40df0855b4a7-recall-and-next-depth}

\begin{itemize}
\tightlist
\item
  Compare a JVP and VJP by naming the direction propagated and the Jacobian multiplication performed.
\item
  Why does a correct gradient still fail to guarantee a global minimum?
\end{itemize}

\section{Source and navigation}\label{appendix-40df0855b4a7-source-and-navigation}

Teaching narrative: \href{Continuous\%20Notes/University\%20-\%20Worlds,\%20Images\%20and\%20Inference\%20-\%20Graphics,\%20Vision\%20and\%20Photography\#Learning\%20a\%20visual\%20classifier\%20from\%20examples}{open the main explanation}. Browse the \href{../02-Foundations-II/main.pdf\#nameddest=appendix-a6d556a4ef50-university-graphics-vision-and-photography---independent-series}{University graphics and vision index (02-Foundations-II, appendix-a6d556a4ef50-university-graphics-vision-and-photography-\/-\/-independent-series)}. Additional chapter connections: \href{Continuous\%20Notes/University\%20-\%20Worlds,\%20Images\%20and\%20Inference\%20-\%20Graphics,\%20Vision\%20and\%20Photography\#Optimization,\%20constraints,\%20and\%20inverse\%20kinematics}{Optimization, constraints, and inverse kinematics}; \href{Continuous\%20Notes/Computing\%20Foundations\%20-\%20From\%20Bits\%20to\%20C++\%20and\%20Systems\%20APIs\#QR\%20turns\%20least\%20squares\%20into\%20a\%20triangular\%20solve\%20plus\%20an\%20irreducible\%20residual}{QR turns least squares into a triangular solve plus an irreducible residual}; \href{Continuous\%20Notes/Computing\%20Foundations\%20-\%20From\%20Bits\%20to\%20C++\%20and\%20Systems\%20APIs\#A\%20compiled\%20interface\%20rather\%20than\%20repeated\%20text}{A compiled interface rather than repeated text}.

Caption/source records: \href{Continuous\%20Notes/University\%20-\%20Worlds,\%20Images\%20and\%20Inference\%20-\%20Graphics,\%20Vision\%20and\%20Photography\#SC-02\%20-\%20Stanford\%20CS231N\%20Spring\%202025\%20Lecture\%202:\%20Image\%20Classification\%20with\%20Linear\%20Classifiers}{SC-02}, \href{Continuous\%20Notes/University\%20-\%20Worlds,\%20Images\%20and\%20Inference\%20-\%20Graphics,\%20Vision\%20and\%20Photography\#SC-03\%20-\%20Stanford\%20CS231N\%20Spring\%202025\%20Lecture\%203:\%20Regularization\%20and\%20Optimization}{SC-03}, \href{Continuous\%20Notes/University\%20-\%20Worlds,\%20Images\%20and\%20Inference\%20-\%20Graphics,\%20Vision\%20and\%20Photography\#SC-04\%20-\%20Stanford\%20CS231N\%20Spring\%202025\%20Lecture\%204:\%20Neural\%20Networks\%20and\%20Backpropagation}{SC-04}.
